To address this problem, a work-orchestration system must satisfy requirements that are not well covered by chat-only or tool-first abstractions. First, it must preserve one durable work thread from demand intake to completion. A buyer should not lose continuity when the workflow moves from clarification to routing, from routing to approval, or from approval to execution. Second, it must represent commercial boundaries explicitly. Approval, pricing, and payment are not side messages. They are part of the work lifecycle and must be modeled as first-class state.

Third, the system must support mixed participation. Real requests may involve human specialists, agents, tools, organizations, and local runtimes in one fulfillment path. The model therefore needs a durable way to represent both who owns the request and which supply-side capabilities respond to it. Fourth, the system must preserve auditable history. Current summary state is useful for product surfaces, but the durable record of what happened, who changed it, and which proof artifacts were attached should remain reconstructible.

Fifth, the system must support proof-bearing completion rather than generated closure. Completion should depend on evidence appropriate to the work mode. For digital work, this may involve code, documents, URLs, or execution traces. For embodied work, it may involve photographs, signatures, geostamped check-ins, access logs, receipts, or witness-confirmed handoffs. The model must therefore represent proof obligations without forcing all proof data inline onto the root object.

Sixth, the system must represent embodied constraints explicitly. These include location requirements, travel or geography constraints, time windows, access dependencies, safety constraints, trust requirements, and the need for human presence. If those requirements exist only in freeform text or remain implicit in the buyer's message, the planner will be tempted to optimize them away. The representation does not have to introduce a new root object, but it does need durable fields or derived structures that allow the system to reason about embodied execution honestly.

Finally, the system must resist premature decomposition. If decomposition happens before a credible lead route is known, the planner is likely to produce generic task trees or tool-shaped work breakdowns. A realistic orchestration system should establish the primary execution path first, then derive bounded sub-work around that path. This motivates the lead-first rule that we use later in the paper.
